\section{Evaluation Protocol and Supplementary Evidence}

The appendix reports protocol checks, auxiliary baseline-transfer evidence, runtime-generated workflow traces, and supplementary broad-benchmark analyses.

\paragraph{Appendix contents.}
This appendix is organized as follows:
\begin{itemize}[leftmargin=12pt, nosep]
    \item Appendix~\ref{subsec:appendix_protocol_fidelity} reports dataset sampling, tool regimes, Wilson guidance, framework comparison, ReAct calibration, backend checks, and protocol-alignment tables.
    \item Appendix~\ref{subsec:appendix_prompt_traj_examples} gives retained prompt and trajectory excerpts for the runtime-workflow setting.
    \item Appendix~\ref{subsec:appendix_process_diagnostics} reports illustrative process signals and separates observed trace evidence from mechanisms that require ablation.
    \item Appendix~\ref{subsec:appendix_runtime_evidence} reports runtime-generated workflow trace statistics and same-task visual evidence.
    \item Appendix~\ref{subsec:appendix_responsible_artifact_use} summarizes artifact use, release planning, and responsible-use notes.
    \item Appendix~\ref{subsec:appendix_broad_benchmark_analyses} provides supplementary broad-benchmark visual analyses.
\end{itemize}

\subsection{Protocol and Fidelity Checks}
\label{subsec:appendix_protocol_fidelity}

Table~\ref{tab:appendix_dataset_sampling} reports the evaluation size used for each dataset. For any source split with more than 400 instances, we evaluate a fixed-random-seed subset of 400 instances. For source splits with at most 400 instances, we evaluate the full split. The full-split evaluations below 400 instances are AIME (30), HumanEval (131), MBPP (341), and GAIA validation (165). The benchmark-balanced average accuracy in Table~\ref{tab:main_exp1_overall} is computed over benchmark-level scores rather than by pooling all evaluated instances.

\begin{table*}[t]
    \centering
    \caption{Dataset sampling protocol for the reported experiments. Splits larger than 400 instances are evaluated on a fixed-random-seed subset of 400 examples; splits with at most 400 instances are evaluated in full.}
    \label{tab:appendix_dataset_sampling}
    \footnotesize
    \setlength{\tabcolsep}{5pt}
    \renewcommand{\arraystretch}{1.08}
    \resizebox{\textwidth}{!}{
    \begin{tabular}{>{\raggedright\arraybackslash}p{0.26\textwidth} >{\raggedright\arraybackslash}p{0.22\textwidth} >{\raggedright\arraybackslash}p{0.22\textwidth} >{\raggedright\arraybackslash}p{0.22\textwidth}}
        \toprule
        \textbf{Dataset} & \textbf{Experiment} & \textbf{Evaluated instances} & \textbf{Sampling rule} \\
        \midrule
        MATH & Broad benchmark & 400 & Fixed-seed subset \\
        AIME & Broad benchmark & 30 & Full split \\
        GSM8K & Broad benchmark & 400 & Fixed-seed subset \\
        DROP & Broad benchmark & 400 & Full split \\
        BBH & Broad benchmark & 400 & Fixed-seed subset \\
        MMLU-Pro & Broad benchmark & 400 & Full split \\
        HumanEval & Broad benchmark & 131 & Full split \\
        MBPP & Broad benchmark & 341 & Full split \\
        HotpotQA & Broad benchmark & 400 & Fixed-seed subset \\
        IFEval & Broad benchmark & 400 & Fixed-seed subset \\
        GAIA validation & Runtime-workflow study & 165 & Full split; levels $(N_1,N_2,N_3)=(53,86,26)$ \\
        \bottomrule
    \end{tabular}
    }
\end{table*}

\begin{table*}[t]
    \centering
    \caption{Benchmark-specific tool regimes used in the reported experiments. Final-answer emission is shared across all runs through the evaluator/answer interface and is omitted from the regime labels.}
    \label{tab:appendix_tool_regimes}
    \footnotesize
    \setlength{\tabcolsep}{5pt}
    \renewcommand{\arraystretch}{1.08}
    \begin{tabular}{>{\raggedright\arraybackslash}p{0.40\textwidth} >{\raggedright\arraybackslash}p{0.18\textwidth} >{\raggedright\arraybackslash}p{0.34\textwidth}}
        \toprule
        \textbf{Dataset(s)} & \textbf{Experiment} & \textbf{Tool regime} \\
        \midrule
        MATH, AIME, GSM8K, DROP, BBH, MMLU-Pro, HumanEval, MBPP, IFEval & Broad benchmark & \texttt{python\_interpreter} only \\
        HotpotQA & Broad benchmark & Expanded full BenchAgent tool registry \\
        GAIA validation (BenchAgent-based systems) & Runtime-workflow study & Expanded full BenchAgent tool registry \\
        GAIA validation (CC-workflow) & Runtime-workflow study & PAE-aligned GAIA-relevant full-capability tool surface \\
        \bottomrule
    \end{tabular}
\end{table*}

\paragraph{Global model-call configuration.}
Unless otherwise noted, all experiments use \texttt{MAX\_TOKEN\_SIZE}=8192, \texttt{temperature}=0.2, and \texttt{top\_p}=1.0.

\begin{table*}[t]
    \centering
    \caption{Approximate Wilson 95\% confidence-interval half-widths for representative evaluated split sizes. Values are maximum half-widths over all Bernoulli accuracies and are used only as interpretation guidance for one-run pass@1 results, not as rerun variance estimates.}
    \label{tab:appendix_wilson_widths}
    \footnotesize
    \setlength{\tabcolsep}{6pt}
    \renewcommand{\arraystretch}{1.08}
    \resizebox{\textwidth}{!}{
    \begin{tabular}{>{\raggedright\arraybackslash}p{0.55\textwidth}cc}
        \toprule
        \textbf{Split / Setting} & \textbf{Evaluated instances} & \textbf{Max Wilson half-width} \\
        \midrule
        Broad benchmarks evaluated on 400-instance fixed-seed subsets & 400 & $\pm$4.9 pt \\
        AIME & 30 & $\pm$16.8 pt \\
        HumanEval & 131 & $\pm$8.4 pt \\
        MBPP & 341 & $\pm$5.3 pt \\
        GAIA Level 1 & 53 & $\pm$13.0 pt \\
        GAIA Level 2 & 86 & $\pm$10.3 pt \\
        GAIA Level 3 & 26 & $\pm$17.9 pt \\
        \bottomrule
    \end{tabular}
    }
\end{table*}

\begin{table*}[t]
    \centering
    \caption{Feature-level comparison with adjacent agent-evaluation infrastructure. The table is qualitative because the systems target different layers: AgentBench and GAIA-style benchmarks define tasks and interaction settings, Harbor emphasizes sandboxed agent/model evaluation, and BenchAgent emphasizes cross-paradigm workflow alignment over existing benchmarks.}
    \label{tab:appendix_framework_feature_comparison}
    \footnotesize
    \setlength{\tabcolsep}{5pt}
    \renewcommand{\arraystretch}{1.08}
    \resizebox{\textwidth}{!}{
    \begin{tabular}{>{\raggedright\arraybackslash}p{0.30\textwidth} c c c}
        \toprule
        \textbf{Feature} & \textbf{AgentBench / task benchmarks} & \textbf{Harbor} & \textbf{BenchAgent} \\
        \midrule
        Existing benchmark reuse rather than new task split & partial & partial & yes \\
        Shared answer contract across compared workflows & partial & partial & yes \\
        Cross-paradigm runtime coverage (single-agent, fixed MAS, evolving MAS) & limited & partial & yes \\
        Token and latency accounting under one protocol & partial & yes & yes \\
        Structured trajectory serialization for workflow analysis & partial & partial & yes \\
        Workflow-topology summaries (agents, edges, depth) & no / limited & no / limited & yes \\
        Tool-surface alignment for coordination-policy comparison & partial & partial & yes \\
        \bottomrule
    \end{tabular}
    }
\end{table*}

Table~\ref{tab:appendix_react_calibration} is a controller-strength calibration, not an alternative denominator for workflow lift. It changes the controller relative to BenchAgent Core, so the scores contextualize the matched anchor used in Table~\ref{tab:main_exp1_overall} rather than replacing the SI comparison, whose target is the effect of changing workflow organization while holding the BenchAgent execution substrate fixed.

\begin{table*}[t]
    \centering
    \caption{Auxiliary controller calibration against a standalone vanilla ReAct controller under aligned BenchAgent tool regimes. MATH, BBH, and HumanEval use the \texttt{python\_interpreter}-only regime, whereas HotpotQA uses the expanded full-tool regime. All rows use the same GPT-4.1 backend and evaluator pipeline as the corresponding BenchAgent anchor.}
    \label{tab:appendix_react_calibration}
    \footnotesize
    \setlength{\tabcolsep}{5pt}
    \renewcommand{\arraystretch}{1.08}
    \resizebox{\textwidth}{!}{
    \begin{tabular}{lcccccc}
        \toprule
        \textbf{Dataset} & \textbf{Evaluated instances} & \textbf{BenchAgent anchor acc.} & \textbf{vanilla ReAct acc.} & \textbf{$\Delta$ acc.} & \textbf{ReAct token / successful case} & \textbf{ReAct time/inst.} \\
        \midrule
        MATH & 400 & 66.75{\scriptsize $\pm$4.6}\% & 78.28{\scriptsize $\pm$4.1}\% & +11.53 pt & 9208.81 & 9.69s \\
        BBH & 400 & 78.25{\scriptsize $\pm$4.0}\% & 85.03{\scriptsize $\pm$3.6}\% & +6.78 pt & 4212.89 & 32.12s \\
        HumanEval & 131 & 84.73{\scriptsize $\pm$6.2}\% & 77.95{\scriptsize $\pm$7.2}\% & -6.78 pt & 59112.01 & 93.98s \\
        HotpotQA & 400 & 62.25{\scriptsize $\pm$4.7}\% & 67.34{\scriptsize $\pm$4.6}\% & +5.09 pt & 8136.38 & 24.19s \\
        \bottomrule
    \end{tabular}
    }
\end{table*}

Tables~\ref{tab:appendix_qwen32b_gaia_generalization} and~\ref{tab:appendix_glm5_gaia_generalization} report GAIA-only backend checks for BenchAgent-compatible workflows. The Qwen3-32B table uses the main-text GAIA metrics to compare the BenchAgent anchor with all fixed and evolving MAS rows inside the same BenchAgent protocol under the global model-call configuration. The GLM-5 table gives a focused BenchAgent backend-strength calibration on GAIA only. The external CC-workflow is intentionally excluded from these backend-swap checks because rerouting that controller to a different backend would change the PAE runtime configuration rather than produce a controlled BenchAgent-internal backend ablation.

\begin{table*}[t]
    \centering
    \caption{\small
        \textbf{GAIA validation pass@1 backend check with Qwen3-32B for BenchAgent-compatible workflows.} The backend check follows the global model-call configuration and reports the same metrics as Table~\ref{tab:main_exp2_gaia}. Tok. reports instance-level average end-to-end token usage, and Time reports instance-level average wall-clock time.
    }
    \label{tab:appendix_qwen32b_gaia_generalization}
    \setlength{\tabcolsep}{5pt}
    \small
    \renewcommand{\arraystretch}{1.2}
    \resizebox{\textwidth}{!}{
        \begin{tabular}{l l cccccc}
            \toprule
            \textbf{Paradigm}          & \textbf{Method}    & \textbf{GAIA-L1}               & \textbf{GAIA-L2} & \textbf{GAIA-L3} & \textbf{Avg./Overall} & \textbf{Tok.} & \textbf{Time} \\
            \midrule
            Single-Agent               & BenchAgent & 47.17{\scriptsize $\pm$13.0}\% & 24.42{\scriptsize $\pm$9.0}\% & 19.23{\scriptsize $\pm$14.7}\% & 30.91{\scriptsize $\pm$7.0}\% & 269178.27 & 704.14s \\
            \midrule
            Evolving MAS               & EvoAgent & 26.42{\scriptsize $\pm$11.6}\% & 23.26{\scriptsize $\pm$8.8}\% & 23.08{\scriptsize $\pm$15.5}\% & 24.24{\scriptsize $\pm$6.5}\% & 925891.35 & 424.01s \\
            Fixed MAS                  & LLM-Debate & 32.08{\scriptsize $\pm$12.2}\% & 12.79{\scriptsize $\pm$7.1}\% & 3.85{\scriptsize $\pm$9.1}\% & 17.58{\scriptsize $\pm$5.8}\% & 265544.08 & 80.60s \\
            Fixed MAS                  & Camel & 39.62{\scriptsize $\pm$12.7}\% & 23.26{\scriptsize $\pm$8.8}\% & 3.85{\scriptsize $\pm$9.1}\% & 25.45{\scriptsize $\pm$6.6}\% & 93501.17 & 389.01s \\
            Fixed MAS                  & AutoGen & 22.64{\scriptsize $\pm$11.0}\% & 16.28{\scriptsize $\pm$7.8}\% & 19.23{\scriptsize $\pm$14.7}\% & 18.79{\scriptsize $\pm$5.9}\% & 310282.87 & 88.51s \\
            Fixed MAS                  & Jarvis & 39.62{\scriptsize $\pm$12.7}\% & 18.60{\scriptsize $\pm$8.2}\% & 23.08{\scriptsize $\pm$15.5}\% & 26.06{\scriptsize $\pm$6.6}\% & 213783.77 & 507.48s \\
            Fixed MAS                  & ChatEval & 28.30{\scriptsize $\pm$11.8}\% & 17.44{\scriptsize $\pm$8.0}\% & 7.69{\scriptsize $\pm$11.0}\% & 19.39{\scriptsize $\pm$6.0}\% & 479251.38 & 207.24s \\
            \bottomrule
        \end{tabular}
    }
    \renewcommand{\arraystretch}{1.0}
\end{table*}

\begin{table*}[t]
    \centering
    \caption{\small
        \textbf{BenchAgent GAIA backend-strength calibration with GLM-5.} The table isolates the single-agent BenchAgent anchor under the GAIA protocol. Level cells report accuracy, instance-level average token usage, and instance-level average wall-clock time.
    }
    \label{tab:appendix_glm5_gaia_generalization}
    \setlength{\tabcolsep}{5pt}
    \small
    \renewcommand{\arraystretch}{1.2}
    \resizebox{\textwidth}{!}{
        \begin{tabular}{l l ccc}
            \toprule
            \textbf{Paradigm} & \textbf{Method} & \textbf{GAIA-L1} & \textbf{GAIA-L2} & \textbf{GAIA-L3} \\
            \midrule
            Single-Agent & BenchAgent & 81.13{\scriptsize $\pm$10.4}\% / 268743.54 / 531.52s & 64.71{\scriptsize $\pm$10.0}\% / 418454.21 / 654.70s & 40.00{\scriptsize $\pm$17.6}\% / 3375974 / 810.14s \\
            \bottomrule
        \end{tabular}
    }
    \renewcommand{\arraystretch}{1.0}
\end{table*}

The backend checks test whether the main workflow-level pattern is specific to one GPT-4.1 setting. They do not introduce a separate conclusion; instead, they provide scoped evidence for the same interpretation reported in the main experiments: MAS lift is not uniformly positive and depends on the interaction among backend reliability, task difficulty, and workflow structure. Under Qwen3-32B, the evaluated MAS rows all trail the BenchAgent anchor overall. Jarvis, Camel, and EvoAgent are closer to the anchor, but Camel's Level~3 drop and the larger gaps for LLM-Debate, AutoGen, and ChatEval are consistent with the view that added coordination can be sensitive to harder multi-step evidence-control tasks. The GLM-5 BenchAgent calibration gives the complementary single-agent reference point, showing that changing the backend can substantially change absolute GAIA level scores while leaving the workflow-comparison question intact. Together, these checks support a bounded conclusion rather than a new claim: agent count or MAS structure alone is not a reliable source of improvement under the tested protocols.

Table~\ref{tab:appendix_gaia_protocol_summary} lists the retained protocol evidence for the GAIA comparison. Fields that cannot be certified from retained artifacts are marked explicitly rather than silently normalized.

\begin{table*}[t]
    \centering
    \caption{Workflow-category terminology used in the paper. The main text uses the fixed and evolving categories consistently.}
    \label{tab:appendix_workflow_taxonomy}
    \footnotesize
    \setlength{\tabcolsep}{4pt}
    \renewcommand{\arraystretch}{1.1}
    \begin{tabular}{>{\raggedright\arraybackslash}p{0.20\textwidth} >{\raggedright\arraybackslash}p{0.43\textwidth} >{\raggedright\arraybackslash}p{0.29\textwidth}}
        \toprule
        \textbf{Category} & \textbf{Definition} & \textbf{Systems in this paper} \\
        \midrule
        Single-agent & One persistent controller maintains context, calls tools, and emits the final answer. & BenchAgent Core \\
        Fixed MAS & Roles, topology, and communication procedure are specified before execution. & CAMEL, AutoGen instantiation, Jarvis/HuggingGPT-style wrapper, ChatEval, LLM-Debate \\
        Evolving MAS & The system searches, mutates, or selects variants within a predefined workflow family. & EvoAgent \\
        Runtime-generated workflow & The controller can create task-specific agents, state artifacts, tool scopes, and verification branches during execution. & CC-workflow with a GPT-4.1 backend \\
        \bottomrule
    \end{tabular}
\end{table*}

\begin{table*}[t]
    \centering
    \caption{Artifact-grounded protocol summary for the GAIA comparison.}
    \label{tab:appendix_gaia_protocol_summary}
    \footnotesize
    \setlength{\tabcolsep}{3.5pt}
    \renewcommand{\arraystretch}{1.04}
    \begin{tabular}{>{\raggedright\arraybackslash}p{0.23\textwidth} >{\raggedright\arraybackslash}p{0.35\textwidth} >{\raggedright\arraybackslash}p{0.35\textwidth}}
        \toprule
        \textbf{Field} & \textbf{BenchAgent-Based Systems} & \textbf{Claude Code Configuration} \\
        \midrule
        Backbone and provider & Official OpenAI \texttt{gpt-4.1-2025-04-14} snapshot via the provider family used in the main experiments & Official OpenAI \texttt{gpt-4.1-2025-04-14} snapshot routed through the same LiteLLM-compatible proxy family; model names appearing only in setup examples are not treated as experiment settings \\
        Model-call parameters & Global \texttt{MAX\_TOKEN\_SIZE}=8192; \texttt{temperature}=0.2; \texttt{top\_p}=1.0 & Global \texttt{MAX\_TOKEN\_SIZE}=8192; retained records show \texttt{temperature}=0.2 and \texttt{top\_p}=1.0 for OpenAI-compatible calls \\
        Invocation and version & BenchAgent runner scripts under the shared evaluation harness & Claude Code VS Code plugin, version 2.1.68; release date is not certified from retained artifacts \\
        Reporting protocol & pass@1, one end-to-end run per instance & pass@1, one end-to-end run per instance \\
        Evaluation split & GAIA validation, Levels 1--3; $(N_1,N_2,N_3)=(53,86,26)$ & GAIA validation, Levels 1--3; $(N_1,N_2,N_3)=(53,86,26)$ \\
        Execution limits & manager \texttt{max\_steps}=15; search \texttt{max\_steps}=10; 600s API timeout; 300s default MCP/server timeout & retained artifacts show prompt-level batching and a 60s shell-command timeout; no uniform controller-level turn, agent-count, wall-clock, or provider-timeout cap is certified, so the PAE comparison is not a resource-budget-matched ablation \\
        Tool and server configuration & expanded full BenchAgent tool registry for GAIA & Claude Code built-in tools plus Agent Teams orchestration; no MCP or external servers were enabled \\
        Permission and sandbox surface & task-level sandbox, tool allowlists, and MCP/server timeouts controlled by BenchAgent & documented VS Code permission profile with interactive prompts bypassed for unattended execution; prompt and settings constrain intended writes to \texttt{outputs/**} and \texttt{logs/**} and deny secrets or destructive deletion \\
        Retry and stopping & model-call retry with \texttt{max\_retries}=3; possible suggestion-guided rerun after failed checks; termination on final answer, step budget, or runtime error & no explicit provider auto-retry field is observed; repeated/resumed calls are counted when present; final extraction uses answer-only \texttt{outputs/task\_id.txt} files, with existing outputs skipped \\
        Token accounting & token usage is averaged over all attempted instances from recorded prompt/completion metadata aggregated across agent messages & usage records expose prompt/completion metadata for main, subagent, tool-emitting, and auxiliary calls; visible instructions, prompts, tool schemas, tool calls, and returned tool results are counted when model-visible \\
        Process evidence & shared BenchAgent execution substrate and structured traces & four of nine retained dumps preserve explicit team-creation, agent, and send-message events, role-scoped subagents, and verifier-stage release; per-subagent internal context is summarized through retained model-call and message records rather than fully replayable \\
        \bottomrule
    \end{tabular}
\end{table*}

\Needspace{7\baselineskip}
\paragraph{Claude Code configuration notes.}
The retained setup artifacts enable experimental Agent Teams and define four named roles: Planner, Solver, Verifier, and Writer. The prompt forbids modifying project code, deleting files, or running destructive/network commands, and requires the Writer to emit final answers and logs under \texttt{outputs/} and \texttt{logs/}. No MCP or external servers are recorded for these runs. Local proxy credentials and API tokens are deliberately omitted from this paper.

\begin{table*}[t]
\centering
\caption{GAIA-critical capability and tool-surface alignment for the main comparison. BenchAgent retained counts are reconstructed from \texttt{log.rar} and the selected Jarvis/ChatEval response-log archive; fuller method-level counts appear in Tables~\ref{tab:appendix_benchagent_gaia_log_counts} and~\ref{tab:appendix_selected_gaia_response_log_stats}.}
\label{tab:appendix_gaia_tool_alignment}
\setlength{\tabcolsep}{3pt}
\renewcommand{\arraystretch}{1.08}
\scriptsize
\resizebox{\textwidth}{!}{
\begin{tabular}{>{\raggedright\arraybackslash}p{0.17\textwidth} >{\raggedright\arraybackslash}p{0.25\textwidth} >{\raggedright\arraybackslash}p{0.27\textwidth} >{\raggedright\arraybackslash}p{0.25\textwidth}}
\toprule
\textbf{Capability} & \textbf{BenchAgent Configured Surface} & \textbf{BenchAgent Retained Evidence} & \textbf{Claude Code Evidence} \\
\midrule
Web search & Web-search and retrieval tools in the GAIA registry. & \texttt{log.rar}: 11{,}050 \texttt{Request URL} events (17.40/inst.). Selected ChatEval: 1{,}836 retrieval-step calls and 3{,}544 executed retrieval records. & Schema exposes \texttt{WebSearch}; retained structured events include 3 calls. \\
Web page fetching / browsing & Web, crawler, browser, and page-reader tools. & Selected ChatEval: 3{,}916 executed external-tool records (24.02/inst.); browser/crawl/archive tools appear among top executed tools. & Web-fetch style capability is visible; parser does not reconstruct every page event. \\
Local file discovery & Benchmark-visible files and file readers. & Discovery is not separated from benchmark file injection in retained BenchAgent logs. & 2 retained \texttt{Glob} calls. \\
Local file reading & Text, Markdown, document, and file readers. & Selected ChatEval: 168 executed file-tool records (1.03/inst.); Jarvis and \texttt{log.rar} do not serialize typed file reads consistently. & 29 retained \texttt{Read} calls; additional inspection often goes through \texttt{Bash}. \\
Shell and code execution & Python/interpreter-style execution through the shared runtime. & \texttt{log.rar}: 159 \texttt{Code Agent Step} and 59 \texttt{Executing code} markers. Selected ChatEval: 2{,}787 structured code-tool calls. & 207 retained \texttt{Bash} events. \\
Structured tables & CSV and sheet-style extraction tools. & CSV/sheet inspection is included in selected ChatEval file-tool records; complete per-format counts are not certified. & Structured files are handled through file tools, shell commands, or Python; no specialized spreadsheet parser is certified. \\
Document / text parsing & Text and document parsing paths. & Text extraction and Markdown conversion appear among selected ChatEval top tools; legacy logs do not replay parser internals. & Direct file reads and shell-based parsing are visible; parser internals are not fully replayable. \\
Writing artifacts & Answer, trajectory, and usage serialization. & Selected response files cover all 328 Jarvis/ChatEval instances and serialize messages, usage, and final answers; \texttt{tool\_calls} are absent. & 65 retained \texttt{Write} events; prompt restricts writes to \texttt{outputs/} and \texttt{logs/}. \\
Team orchestration & Fixed, dynamic, or debate-style MAS wrappers. & Selected responses: ChatEval has 1{,}878 model calls (11.52/inst.); Jarvis has 165 (1.00/inst.). & 7 \texttt{TeamCreate}, 31 \texttt{Agent}, and 35 \texttt{SendMessage} events. \\
Task/state management & Manager state, step budgets, messages, tool calls, and evaluator outputs. & Selected responses: 3{,}756 ChatEval messages (23.04/inst.) and 330 Jarvis messages (2.00/inst.). & Role-scoped messages, written artifacts, final-answer files, and \texttt{TodoWrite} schema are visible. \\
External servers & Optional MCP/server-backed registered tools. & Retained logs do not separate MCP/server-backed calls from ordinary registered tool calls. & No MCP or external servers were enabled. \\
\bottomrule
\end{tabular}
}
\renewcommand{\arraystretch}{1.0}
\end{table*}

\begin{table*}[t]
    \centering
    \caption{GAIA \texttt{log.rar} structure and retained event counts for BenchAgent-based systems. Each cell reports total count with mean per retained instance in parentheses. These legacy logs do not serialize structured \texttt{STEP\_LOG} or \texttt{TOOL\_CALL name=...} records, so counts are conservative retained-log evidence rather than complete hidden-controller tool traces. ``Code-agent steps'' counts logged \texttt{Code Agent Step} markers; ``Direct external'' counts \texttt{Calling tool: ...} markers after excluding \texttt{final\_answer}; ``Wiki/API req.'' counts logged \texttt{Request URL:} events; ``Code exec.'' counts \texttt{Executing code:}; and ``Tool errors'' counts logged tool-execution errors. For Jarvis and ChatEval, Table~\ref{tab:appendix_selected_gaia_response_log_stats} separately reports the selected-batch response/log parser output used for the GAIA table.}
    \label{tab:appendix_benchagent_gaia_log_counts}
    \footnotesize
    \setlength{\tabcolsep}{4pt}
    \renewcommand{\arraystretch}{1.08}
    \resizebox{\textwidth}{!}{
    \begin{tabular}{lcccccc}
        \toprule
        \textbf{Method} & \textbf{Retained instances} & \textbf{Code-agent steps} & \textbf{Direct external} & \textbf{Wiki/API req.} & \textbf{Code exec.} & \textbf{Tool errors} \\
        \midrule
        Single Agent & 127 & 45 (0.35) & 116 (0.91) & 2301 (18.12) & 33 (0.26) & 7 (0.06) \\
        EvoAgent & 127 & 50 (0.39) & 0 (0.00) & 2429 (19.13) & 0 (0.00) & 15 (0.12) \\
        LLM-Debate & 127 & 15 (0.12) & 0 (0.00) & 1987 (15.65) & 2 (0.02) & 9 (0.07) \\
        Camel & 127 & 21 (0.17) & 10 (0.08) & 2433 (19.16) & 7 (0.06) & 7 (0.06) \\
        AutoGen & 127 & 28 (0.22) & 18 (0.14) & 1900 (14.96) & 17 (0.13) & 10 (0.08) \\
        \bottomrule
    \end{tabular}
    }
    \renewcommand{\arraystretch}{1.0}
\end{table*}

\begin{table*}[t]
    \centering
    \caption{Selected GAIA response/log statistics for Jarvis and ChatEval. The six selected result JSON files in the supplementary archive define the evaluated batches; matching response files cover all 165 Jarvis and 163 ChatEval retained instances. Response files reliably serialize messages, model-call usage, and token usage, but they do not persist \texttt{tool\_calls}. Tool-use evidence is therefore extracted from logs. ``Step external'' counts structured \texttt{STEP\_LOG} tool calls excluding \texttt{final\_answer}; ``executed external'' counts tool-manager \texttt{TOOL\_CALL} records excluding \texttt{final\_answer}; and ``direct/code markers'' counts unstructured \texttt{Calling tool: ...} lines excluding \texttt{final\_answer} plus \texttt{Executing code:} markers. Jarvis logs do not serialize structured \texttt{STEP\_LOG} or \texttt{TOOL\_CALL} records for these batches, so only unstructured markers are available for Jarvis.}
    \label{tab:appendix_selected_gaia_response_log_stats}
    \footnotesize
    \setlength{\tabcolsep}{4pt}
    \renewcommand{\arraystretch}{1.08}
    \resizebox{\textwidth}{!}{
    \begin{tabular}{lcccccc}
        \toprule
        \textbf{Method} & \textbf{Selected instances} & \textbf{Matched responses} & \textbf{Model calls} & \textbf{Step external} & \textbf{Executed external} & \textbf{Direct/code markers} \\
        \midrule
        Jarvis & 165 & 165 (1.00) & 165 (1.00) & 0 (not serialized) & 0 (not serialized) & 22 (0.13) \\
        ChatEval & 163 & 163 (1.00) & 1878 (11.52) & 4623 (28.36) & 3916 (24.02) & 1189 (7.29) \\
        \bottomrule
    \end{tabular}
    }
    \renewcommand{\arraystretch}{1.0}
\end{table*}

\begin{table*}[t]
    \centering
    \caption{Retained Claude Code GAIA prompt and output contract. The table records the artifact content used for the reported runs; credentials and local proxy details are intentionally excluded.}
    \label{tab:appendix_claude_prompt_contract}
    \footnotesize
    \setlength{\tabcolsep}{4pt}
    \renewcommand{\arraystretch}{1.08}
    \begin{tabular}{>{\raggedright\arraybackslash}p{0.23\textwidth} >{\raggedright\arraybackslash}p{0.69\textwidth}}
        \toprule
        \textbf{Prompt Component} & \textbf{Retained Content / Operational Effect} \\
        \midrule
        Autonomy rule & The prompt instructs the system to operate autonomously as an Agent Team and continue without waiting for user confirmation unless a critical error occurs. \\
        Team roles & Four named roles are specified: Planner, Solver, Verifier, and Writer. Planner inspects the project and task split; Solver reasons over each task; Verifier checks the proposed answer; Writer emits final answer files and logs. \\
        Input-reading rule & The GAIA prompt asks the Planner to inspect JSONL split files using shell-visible file operations such as \texttt{head}, \texttt{sed}, and \texttt{tail}, rather than relying on an implicit hidden reader. \\
        Split and batching & Retained prompts include both a first-10-task batch contract and a GAIA validation contract over \texttt{split\_validation/gaia1\_2.jsonl}; existing output files are skipped rather than overwritten. \\
        Write scope & The prompt restricts writes to \texttt{outputs/} and \texttt{logs/}. It explicitly forbids modifying project code, task data, prompts, source directories, configuration files, or existing benchmark assets. \\
        Command restrictions & Destructive file deletion and network-download commands are forbidden by prompt-level rules; the retained settings also deny secrets and destructive deletion patterns. \\
        Final extraction & Final answers are written as concise answer-only \texttt{outputs/<task\_id>.txt} artifacts, with per-task reasoning or evidence placed under \texttt{logs/}. The user-facing final report lists processed, skipped, and generated outputs. \\
        \bottomrule
    \end{tabular}
    \renewcommand{\arraystretch}{1.0}
\end{table*}

\paragraph{Claude Code retained-log token accounting.}
The retained LiteLLM-style records make the visible token boundary auditable for the PAE run. We therefore report CC-workflow cost as retained-log accounting rather than as provider-independent billing or a strict resource-budget ablation against BenchAgent. The counted usage includes model-visible instructions, the GAIA mission prompt, visible \texttt{system-reminder} context, tool schemas, tool-call JSON, subagent/delegation model calls, auxiliary title-generation calls, and tool results once returned to later prompts. Local shell execution is not tokenized until its output is fed back to the model. In the deduplicated batch, the logs contain 128 subagent model calls, 206 calls that emitted tool calls, 236 calls whose prompts include historical \texttt{tool\_use}/\texttt{tool\_result} context, and 22 title-generation calls. We did not observe an actual compaction-summary event in this batch; invisible provider-side policy remains outside the retained logs and is not used for mechanism attribution.

\subsection{Prompt and Trajectory Excerpts}
\label{subsec:appendix_prompt_traj_examples}

This subsection records compact prompt and trajectory excerpts used to audit the reported runs. Table~\ref{tab:appendix_benchagent_core_prompt} summarizes the BenchAgent Core prompt contract from \texttt{bench/mas\_arena/agents/prompts.yaml} and \texttt{bench/mas\_arena/agents/bench\_agent.py}. Tables~\ref{tab:appendix_reproduced_mas_contracts} and~\ref{tab:appendix_evoagent_contract} extend the record to the reproduced fixed and evolving MAS wrappers. Table~\ref{tab:appendix_retained_trajectory_excerpt} gives three retained single-task trajectory excerpts from the response and structured-log artifacts, and Table~\ref{tab:appendix_additional_trajectory_excerpts} adds further excerpts covering the ChatEval debate structure and the CC-workflow verifier-release gate. Long task text, local paths, and hidden tool schemas are shortened only for readability.

\begin{table*}[t]
    \centering
    \caption{BenchAgent Core prompt and execution contract. These fields are copied or closely paraphrased from the retained prompt templates and runner code, rather than inferred from results.}
    \label{tab:appendix_benchagent_core_prompt}
    \footnotesize
    \setlength{\tabcolsep}{4pt}
    \renewcommand{\arraystretch}{1.08}
    \begin{tabular}{p{0.20\textwidth}p{0.29\textwidth}p{0.43\textwidth}}
        \toprule
        \textbf{Prompt / code field} & \textbf{Retained wording or behavior} & \textbf{Evaluation role} \\
        \midrule
        Task wrapper & ``Answer this question correctly. You have all the tools needed to find the right answer.'' & Gives every benchmark instance the same task-facing wrapper before workflow-specific execution begins. \\
        Tool priority & ``Always prioritize using available tools over writing custom code.'' & Prevents compared systems from silently bypassing the tool surface through ad hoc local code when a registered tool is intended. \\
        File-access rule & ``If the task involves files or attachments, you MUST use appropriate tools to access them.'' & Keeps file-bearing GAIA and document tasks inside the logged tool interface rather than hidden file reads. \\
        Search delegation & The manager can call \texttt{search\_agent("Your detailed request")} for web research and online documents. & Separates manager reasoning from delegated web research while preserving a recorded handoff. \\
        Direct answer contract & The manager must call \texttt{final\_answer("Your answer")} and keep final answers concise. & Normalizes final-answer extraction across benchmarks and workflows. \\
        Runtime limits & The manager uses \texttt{max\_steps}=15 and the search agent uses \texttt{search\_max\_steps}=10 unless configured otherwise. & Makes step budgets explicit for the SI comparison. \\
        Verification pressure & The task wrapper says failures or \texttt{None found} are not acceptable and asks the agent to verify when needed. & Encourages evidence checking without changing the final-answer-only evaluator. \\
        Trace serialization & Returned records include message history, final answer, manager steps, search-agent steps, and usage metadata when available. & Provides the process evidence used in the main text and appendix diagnostics. \\
        \bottomrule
    \end{tabular}
    \renewcommand{\arraystretch}{1.0}
\end{table*}

Tables~\ref{tab:appendix_reproduced_mas_contracts} and~\ref{tab:appendix_evoagent_contract} extend the prompt-contract record to the other compared workflows. The fixed MAS table makes the ``roles and topology specified before execution'' characterization from Table~\ref{tab:appendix_workflow_taxonomy} concrete: each system names its agents, assigns distinct or homogeneous instruction sets, and fixes the communication procedure before any task input is seen. The EvoAgent table documents the base-prompt seeds and the LLM-driven crossover and mutation operations, clarifying that the evolutionary search is confined to the prompt-scaffold space while the BenchAgent execution substrate and tool surface remain fixed across all candidates.

\begin{table*}[t]
    \centering
    \caption{Retained agent role and execution contracts for the four reproduced fixed MAS wrappers. Instructions are copied or closely paraphrased from the retained source code rather than inferred from results. ``LLM-only'' indicates no external tools beyond the \texttt{final\_answer} interface.}
    \label{tab:appendix_reproduced_mas_contracts}
    \footnotesize
    \setlength{\tabcolsep}{3.5pt}
    \renewcommand{\arraystretch}{1.12}
    \resizebox{\textwidth}{!}{
    \begin{tabular}{>{\raggedright\arraybackslash}p{0.10\textwidth} >{\raggedright\arraybackslash}p{0.14\textwidth} >{\raggedright\arraybackslash}p{0.30\textwidth} >{\raggedright\arraybackslash}p{0.26\textwidth} >{\raggedright\arraybackslash}p{0.13\textwidth}}
        \toprule
        \textbf{System} & \textbf{Agent instances} & \textbf{Role-level instruction content} & \textbf{Communication and termination} & \textbf{Tools and limits} \\
        \midrule
        Jarvis & Single core manager & \textbf{(1) Analyze \& Plan:} briefly outline steps before writing code or calling tools. \textbf{(2) Execute:} implement the plan using Python interpreter and search tools proactively. \textbf{(3) Verify \& Summarize:} check results and provide a concise final answer. & BenchAgent single-step execution; no inter-agent handoff; \texttt{final\_answer} tool terminates. & Python + web search; \texttt{max\_steps}=15 \\
        \addlinespace
        ChatEval & Math Expert (tools); Logic Expert (LLM-only); Critical Thinking Expert (LLM-only); ResultExtractor & \textbf{Math Expert:} analyze mathematically; verify with Python; use \texttt{web\_search} for current data; use file tools for attachments. \textbf{Logic Expert:} analyze logical structure and implicit conditions; be extremely concise; no tools. \textbf{Critical Thinking Expert:} analyze from multiple angles; identify traps; be extremely concise; no tools. & Sequential per-agent contributions each round; all agents see the shared debate history; ResultExtractor aggregates the debate into a single final answer. & Math Expert: Python + web search; Logic/Critical: LLM-only \\
        \addlinespace
        LLM-Debate & 2 homogeneous BenchAgents (\texttt{debate\_agent\_1/2}) + Aggregator & Both agents receive the same base instruction: ``You are a helpful AI assistant.'' No persona differentiation. Aggregator synthesizes the debate history. & Multi-round debate (\texttt{rounds\_num}=3); agents alternate; aggregator emits final answer after the last round. & BenchAgent-configured tool surface per agent \\
        \addlinespace
        CAMEL & User role (instruction giver) + Assistant role (task executor) & \textbf{User:} evaluates and gives feedback; terminates with \texttt{TASK\_FINISHED} when satisfied. Retained prompt: ``If you feel the assistant's response is satisfactory, please include `TASK\_FINISHED'.'' \textbf{Assistant:} ``Respond with clarity and accuracy.'' & Iterative user/assistant conversation; exits on \texttt{TASK\_FINISHED} signal; synthesizer extracts answer. & BenchAgent-configured tools for Assistant \\
        \bottomrule
    \end{tabular}
    }
    \renewcommand{\arraystretch}{1.0}
\end{table*}

\begin{table*}[t]
    \centering
    \caption{Retained EvoAgent evolutionary search contract. The search is confined to the prompt-scaffold space; the BenchAgent ReAct core, tool surface, step budget, and \texttt{final\_answer} interface are held fixed across all candidate agents throughout the evolutionary search.}
    \label{tab:appendix_evoagent_contract}
    \footnotesize
    \setlength{\tabcolsep}{4pt}
    \renewcommand{\arraystretch}{1.12}
    \begin{tabular}{>{\raggedright\arraybackslash}p{0.22\textwidth} >{\raggedright\arraybackslash}p{0.70\textwidth}}
        \toprule
        \textbf{Field} & \textbf{Retained content} \\
        \midrule
        Initial population & 3 agents (EVO-1, EVO-2, EVO-3), each seeded from a distinct base-prompt template. \\
        Seed prompt EVO-1 & ``You are a mathematics expert, skilled in solving mathematical problems. Please think step by step and solve the problem.'' \\
        Seed prompt EVO-2 & ``You are a logical reasoning expert, skilled in analyzing problems and finding solutions. Please provide detailed reasoning process.'' \\
        Seed prompt EVO-3 & ``You are a problem-solving expert, skilled in breaking down complex problems into simple steps. Please clearly show your thinking process.'' \\
        Crossover & An LLM call receives both parent system prompts and generates a new prompt that ``combines the strengths'' of each parent; the offspring agent inherits the merged reasoning emphasis. \\
        Mutation & An LLM call receives one parent prompt and is asked to ``create a mutated agent configuration that is different from the parent but still effective''; the output must include a \texttt{system\_prompt} field. \\
        Selection & Agents achieving better task-level performance are retained as parents for subsequent crossover and mutation. \\
        Fixed substrate & BenchAgent ReAct core, registered tools, API settings, \texttt{max\_steps} budget, and \texttt{final\_answer} contract are identical across all candidate agents; only the \texttt{system\_prompt} field varies across the evolutionary population. \\
        \bottomrule
    \end{tabular}
    \renewcommand{\arraystretch}{1.0}
\end{table*}

\begin{table*}[t]
    \centering
    \caption{Three representative retained single-task trajectory excerpts. Rows preserve within-task order and report the artifact-visible event rather than a reconstructed ideal workflow.}
    \label{tab:appendix_retained_trajectory_excerpt}
    \footnotesize
    \setlength{\tabcolsep}{3pt}
    \renewcommand{\arraystretch}{1.08}
    \begin{tabular}{p{0.09\textwidth}p{0.08\textwidth}p{0.17\textwidth}p{0.58\textwidth}}
        \toprule
        \textbf{Trace} & \textbf{Step} & \textbf{Visible actor} & \textbf{Retained single-task event excerpt} \\
        \midrule
        Core-1 & 1 & user wrapper & The retained query asks the Girls Who Code percentage-change question under the shared BenchAgent wrapper: use tools if needed, verify when needed, and return the answer directly. Source artifact: \texttt{bench\_agent\_7d4a...}. \\
        Core-1 & 2 & final-answer event & The final-answer event emits \texttt{22} with recorded usage metadata; the retained result is correct against the ground truth \texttt{22}. \\
        Core-2 & 1 & user wrapper & A text-grid task asks the agent to read all letters left-to-right from a 5-by-7 block and recover the sentence. Source artifact: \texttt{jarvis\_50ad...}. \\
        Core-2 & 2 & final-answer event & The final-answer event emits \texttt{THE SEA GULL GLIDED PEACEFULLY TO MY CHAIR}, matching the expected sentence up to capitalization and spacing. \\
        Runtime-1 & 1 & team-lead / \texttt{TeamCreate} & The controller creates a GAIA validation team for \texttt{split\_validation/level2\_part7.jsonl}, with planner, solver, verifier, and writer roles. Source artifact: \texttt{3.json}. \\
        Runtime-1 & 2 & team-lead / \texttt{Agent} & The planner is spawned with a bash-only JSONL extraction instruction and asked to return task ids and raw question lines. \\
        Runtime-1 & 3 & team-lead / \texttt{Bash} & The retained command \texttt{head -n 3 level2\_part7.jsonl} exposes the source-task read; the first visible task id begins \texttt{a7feb290{\ldots}}. \\
        Runtime-1 & 4 & team-lead / \texttt{SendMessage} & The lead sends the planner a follow-up to return parseable JSON lines and skip tasks with existing \texttt{outputs/<task\_id>.txt}, making the handoff and output policy explicit. \\
        \bottomrule
    \end{tabular}
    \renewcommand{\arraystretch}{1.0}
\end{table*}

Table~\ref{tab:appendix_additional_trajectory_excerpts} extends the trajectory record to the fixed and runtime-generated workflow paradigms in more detail. The ChatEval-1 trace illustrates how the three-expert debate converges toward a single answer across rounds for a GAIA book-title retrieval task, connecting directly to the token and model-call overhead reported in Table~\ref{tab:appendix_selected_gaia_response_log_stats}. The Runtime-2 and Runtime-3 traces come from sub-agent perspectives within the same CC-workflow run (\texttt{6.json}): Runtime-2 records the Solver role reasoning and forwarding answers to the Verifier, while Runtime-3 records the Writer role receiving Verifier-confirmed packets and enforcing the output contract. Together they provide the artifact-level support for the ``verifier-stage control'' mechanism entry in Table~\ref{tab:gaia_mechanism_evidence}.

\begin{table*}[t]
    \centering
    \caption{Additional retained trajectory excerpts for the fixed and runtime-generated workflow paradigms. ChatEval-1 comes from the extracted GAIA debate logs; Runtime-2 and Runtime-3 come from two sub-agent perspectives in the same \texttt{6.json} CC-workflow dump. Rows preserve within-task order; long content, paths, and tool schemas are shortened for readability.}
    \label{tab:appendix_additional_trajectory_excerpts}
    \footnotesize
    \setlength{\tabcolsep}{3pt}
    \renewcommand{\arraystretch}{1.08}
    \begin{tabular}{>{\raggedright\arraybackslash}p{0.11\textwidth}p{0.06\textwidth}p{0.19\textwidth}p{0.56\textwidth}}
        \toprule
        \textbf{Trace} & \textbf{Step} & \textbf{Visible actor} & \textbf{Retained event excerpt} \\
        \midrule
        \mbox{ChatEval-1} & 1 & user wrapper & A GAIA book-title retrieval task is routed to the three-expert debate team under the shared BenchAgent task wrapper. Source artifact: \texttt{chateval\_newcore\_gaia\_search} debug log. \\
        \mbox{ChatEval-1} & 2 & Math Expert / Round 1 & Math Expert performs a web search and retrieves: ``[1] James Beard Award-winning food journalists Jane and Michael Stern published a book titled `500 Things to Eat Before It's Too Late' in which the restaurant is recommended.'' \\
        \mbox{ChatEval-1} & 3 & Logic Expert / Round 1 & Logic Expert emits the minimal answer form ``the book title only,'' relying on the Math Expert's search evidence without issuing a redundant tool call. \\
        \mbox{ChatEval-1} & 4 & Critical Thinking Expert / Round 1 & Critical Thinking Expert emits the normalized form: ``should be: `five hundred things to eat before it's too late.' '' \\
        \mbox{ChatEval-1} & 5 & All agents / Round 2 & All three experts converge on ``five hundred things to eat before it is too late.'' The ResultExtractor aggregates the debate history into a single final-answer event. \\
        \midrule
        \mbox{Runtime-2} & 1 & Solver / role init & Solver role receives its scope: for each incoming task (task\_id, question), solve strictly from the question and attached context, format the answer per GAIA requirements, and send only task\_id, answer, and reasoning to the Verifier. Source artifact: \texttt{6.json} (sub-agent perspective). \\
        \mbox{Runtime-2} & 2 & Solver / teammate-message & Planner forwards task \texttt{48eb8242\ldots}: ``How many nonindigenous crocodiles were found in Florida from the year 2000 through 2020? (USGS Nonindigenous Aquatic Species database.)'' \\
        \mbox{Runtime-2} & 3 & Solver / \texttt{SendMessage} & Solver sends to Verifier: ``Answer: 11; Task ID: \texttt{48eb8242\ldots}; Reasoning: According to the USGS Nonindigenous Aquatic Species database, there were 11 confirmed occurrences of nonindigenous crocodiles in Florida from 2000 through 2020.'' \\
        \mbox{Runtime-2} & 4 & Solver / teammate-message & Planner forwards task \texttt{c8b7e059\ldots}: ``The work referenced in footnote 397 of Federico Lauria's 2014 dissertation is also the source of two Smithsonian paintings. What is the sum of the paintings' years?'' \\
        \mbox{Runtime-2} & 5 & Solver / \texttt{SendMessage} & Solver sends to Verifier: ``Answer: 23; Task ID: \texttt{c8b7e059\ldots}; Reasoning: Footnote 397 of Federico Lauria's 2014 dissertation references Dante's \emph{Inferno}; the two Smithsonian paintings are `The Fifth Circle' and `The City of Dis'; 23 was computed from the associated publication years.'' \\
        \midrule
        \mbox{Runtime-3} & 1 & Writer / role init & Writer role receives team scope: check for existing \texttt{outputs/<task\_id>.txt} before writing; write only verified answer packets forwarded by the Verifier; save full reasoning to \texttt{logs/<task\_id>.log}. Source artifact: \texttt{6.json} (sub-agent perspective). \\
        \mbox{Runtime-3} & 2 & Writer / teammate-message & Verifier forwards confirmed packet for task \texttt{48eb8242\ldots}: ``Answer: 11; Reasoning: According to the USGS Nonindigenous Aquatic Species database, there were 11 confirmed occurrences\ldots'' The message is gated by the Verifier's format and correctness check before reaching the Writer. \\
        \mbox{Runtime-3} & 3 & Writer / \texttt{Read} & Writer checks whether \texttt{outputs/48eb8242\ldots.txt} already exists before writing, enforcing the skip-existing output policy from the prompt contract. \\
        \mbox{Runtime-3} & 4 & Writer / teammate-message & Verifier forwards confirmed packet for task \texttt{c8b7e059\ldots}: ``Answer: 23; Reasoning: Footnote 397 of Federico Lauria's 2014 dissertation references Dante's \emph{Inferno}\ldots'' \\
        \mbox{Runtime-3} & 5 & Writer / \texttt{Write} & Writer writes the verified answer to \texttt{outputs/c8b7e059\ldots.txt}; the Verifier's confirmation is the necessary precondition for any write event, implementing the verifier-release gate visible in the \texttt{6.json} team run statistics. \\
        \bottomrule
    \end{tabular}
    \renewcommand{\arraystretch}{1.0}
\end{table*}

Table~\ref{tab:appendix_baseline_fidelity} reports an auxiliary transfer sanity check for selected MAS baselines. It compares official implementations or officially reported task-level numbers against matched non-BenchAgent re-instantiations on six shared datasets. These numbers are not the BenchAgent results in Table~\ref{tab:main_exp1_overall}; the check is scoped to task-level plausibility because the original reports and our transferred runs do not always share the same underlying base model generation. Its purpose is to verify that the transferred wrappers remain in a plausible task-level range before they are placed under the controlled BenchAgent protocol.

\begin{table*}[t]
    \centering
    \caption{Auxiliary transfer sanity check for selected representative baselines on six shared datasets. Scores are percentages. The independent transfer rows are separate non-BenchAgent checks and are not the Table~\ref{tab:main_exp1_overall} results. ``Mean Abs. Diff.'' and ``Max Diff.'' are absolute percentage-point differences over MATH, AIME, DROP, MMLU-Pro, BBH, and HumanEval; they are range diagnostics for transfer plausibility rather than code-identity estimates.}
    \label{tab:appendix_baseline_fidelity}
    \small
    \setlength{\tabcolsep}{4.5pt}
    \renewcommand{\arraystretch}{1.1}
    \resizebox{\textwidth}{!}{
    \begin{tabular}{llccccccccc}
        \toprule
        \textbf{Method} & \textbf{Setting} & \textbf{MATH} & \textbf{AIME} & \textbf{DROP} & \textbf{MMLU-Pro} & \textbf{BBH} & \textbf{HumanEval} & \textbf{Avg.} & \textbf{Mean Abs. Diff.} & \textbf{Max Diff.} \\
        \midrule
        \multirow{2}{*}{EvoAgent} & Official & 68.00 & 16.67 & 88.00 & 60.00 & 75.00 & 72.28 & 63.32 & \multirow{2}{*}{6.79} & \multirow{2}{*}{12.72} \\
         & Independent transfer & 65.00 & 26.70 & 80.00 & 60.00 & 82.00 & 85.00 & 66.45 &  &  \\
        \midrule
        \multirow{2}{*}{Jarvis} & Official & 59.00 & 6.67 & 85.00 & 55.00 & 55.00 & 92.00 & 58.78 & \multirow{2}{*}{2.89} & \multirow{2}{*}{5.00} \\
         & Independent transfer & 55.00 & 10.00 & 85.00 & 50.00 & 54.00 & 88.00 & 57.00 &  &  \\
        \midrule
        \multirow{2}{*}{LLM-debate} & Official & 53.00 & 13.30 & 84.00 & 66.00 & 79.00 & 88.00 & 63.88 & \multirow{2}{*}{2.90} & \multirow{2}{*}{3.37} \\
         & Independent transfer & 50.00 & 16.67 & 87.00 & 63.00 & 82.00 & 86.00 & 64.11 &  &  \\
        \midrule
        \multirow{2}{*}{AutoGen} & Official & 49.00 & 10.00 & 87.00 & 64.00 & 77.00 & 77.00 & 60.67 & \multirow{2}{*}{5.12} & \multirow{2}{*}{11.00} \\
         & Independent transfer & 52.00 & 16.70 & 91.00 & 59.00 & 66.00 & 78.00 & 60.45 &  &  \\
        \bottomrule
    \end{tabular}
    }
\end{table*}

The independent transfers fall within a small-to-moderate range of the official reference numbers, but this table is only a sanity check. The controlled evidence in the main paper comes from Table~\ref{tab:main_exp1_overall}, where the compared systems are evaluated under the shared BenchAgent execution substrate.

\subsection{Illustrative Process Signals and Mechanism Evidence}
\label{subsec:appendix_process_diagnostics}

Table~\ref{tab:appendix_process_diagnostics} gives process context for broad-benchmark patterns in Table~\ref{tab:main_exp1_overall}. Because the retained broad-benchmark artifacts do not expose a unified per-wrapper process summary, the measured-quantity column reports illustrative counts from retained GAIA traces for the same wrappers. These rows are used to interpret plausible workflow behavior, not to claim that the GAIA counts causally explain the broad-benchmark scores.

\begin{table*}[t]
\centering
\caption{Illustrative process signals for the workflow wrappers in Main Experiment~1. Counts are from retained GAIA traces for the same wrappers under BenchAgent logging (Appendix Tables~\ref{tab:appendix_benchagent_gaia_log_counts} and~\ref{tab:appendix_selected_gaia_response_log_stats}) and should be read as cross-experiment illustrative evidence, not matched broad-benchmark measurements or causal explanations.}
\label{tab:appendix_process_diagnostics}
\footnotesize
\setlength{\tabcolsep}{3.5pt}
\renewcommand{\arraystretch}{1.12}
\begin{tabular}{p{0.22\textwidth}p{0.22\textwidth}p{0.25\textwidth}p{0.23\textwidth}}
\toprule
\textbf{Observed broad-result pattern} & \textbf{Illustrative wrapper-level signal} & \textbf{Comparable trace quantity} & \textbf{Bounded interpretation} \\
\midrule
EvoAgent has the only positive benchmark-balanced workflow lift, but its +1.44-point margin is within one-run uncertainty. & On retained GAIA runs, BenchAgent logs 19.13 Wiki/API requests and 0.12 tool errors per instance (Table~\ref{tab:appendix_benchagent_gaia_log_counts}). & The evolving wrapper varies workflow and prompt scaffolds inside a designer-specified search space. & The broad-suite gain is better read as task-specific scaffold recovery than a general multi-agent advantage. \\
EvoAgent's largest gain is on BBH, while it trails on AIME, GSM8K, and HumanEval. & Retained GAIA logs record 0.39 code-agent steps per instance (Table~\ref{tab:appendix_benchagent_gaia_log_counts}). & Task-level traces separate where the workflow family helps from where it adds overhead. & Workflow search can locate a reasoning scaffold for some tasks but underperform when exact calculation dominates. \\
LLM-Debate leads on HumanEval and MATH but not on most other tasks. & BenchAgent logs 15.65 Wiki/API requests, 0.12 code-agent steps, and 0.02 code executions per instance (Table~\ref{tab:appendix_benchagent_gaia_log_counts}). & Independent proposal and aggregation stages are visible under the same evaluator. & Debate-style redundancy helps when alternatives can be checked locally, but is not a universal substitute for state preservation. \\
ChatEval is strongest on IFEval but has the highest token usage and trails the anchor overall. & Retained GAIA batches show 23.04 messages, 11.52 model calls, and 24.02 external-tool records per instance (Table~\ref{tab:appendix_selected_gaia_response_log_stats}). & Repeated discussion expands the message history under shared accounting. & Coordination can buy instruction-following robustness in one slice while erasing workflow lift overall. \\
Jarvis is faster and lighter but still loses accuracy. & Retained GAIA batches show 2.00 messages, 1.00 model call, and 0.13 direct/code markers per instance (Table~\ref{tab:appendix_selected_gaia_response_log_stats}). & Its predefined handoff exposes low coordination overhead. & Shallow coordination can be efficient but does not prevent missing task-specific reasoning modes. \\
Camel is faster and lighter but still loses accuracy. & BenchAgent logs 19.16 Wiki/API requests, 0.08 direct external-tool markers, and 0.06 code executions per instance (Table~\ref{tab:appendix_benchagent_gaia_log_counts}). & A low-overhead fixed-handoff pattern without the long discussion history of ChatEval. & Fixed handoff saves tokens but does not guarantee that the right evidence reaches the final decision stage. \\
\bottomrule
\end{tabular}
\renewcommand{\arraystretch}{1.0}
\end{table*}

Table~\ref{tab:gaia_mechanism_evidence} separates observed trace signals from causal claims that require controlled ablation for the CC-workflow GAIA result.

\begin{table*}[t]
\centering
\caption{Evidence status for candidate mechanisms behind the CC-workflow GAIA result, with concrete ablation paths needed to resolve each claim.}
\label{tab:gaia_mechanism_evidence}
\footnotesize
\setlength{\tabcolsep}{2.5pt}
\begin{tabular}{p{0.17\textwidth}p{0.27\textwidth}p{0.24\textwidth}p{0.22\textwidth}}
\toprule
\textbf{Candidate mechanism} & \textbf{Observed trace evidence} & \textbf{Remaining confound} & \textbf{Ablation needed} \\
\midrule
Runtime delegation & Four structured dumps expose \texttt{TeamCreate}, \texttt{Agent}, and \texttt{SendMessage} events with 6.25 subagents on average. & No ablation disables delegation or fixes the agent count. & Compare no-subagent, fixed-agent-count, and adaptive-delegation variants under the same tools and answer contract. \\
Persistent artifacts & Retained traces include \texttt{Read}, \texttt{Write}, and evidence re-opening before finalization in some cases. & File tooling is bundled with the external runtime and task environment. & Disable writes or replace files with in-memory state while holding the solver and tool surface fixed. \\
Verifier-stage control & Two structured dumps show non-idle verifier emissions before final answer release. & The retained traces do not expose every recovery loop or failed verification path. & Ablate the verifier release gate or force final-answer emission before verification and measure error changes. \\
Context management & Some traces preserve task evidence across role boundaries and later checks. & Provider-side context compaction or bookkeeping is not fully observable. & Replay matched evidence packets with controlled context-window and compaction policies. \\
\bottomrule
\end{tabular}
\end{table*}

Table~\ref{tab:appendix_failure_step_distribution} quantifies where errors are first introduced within each wrapper, using the binary-search failure attribution logs from retained GAIA runs. For each failed instance, the binary-search probe identifies the earliest conversation step at which the agent's response diverges from the ground truth. Step~1 corresponds to the first agent turn (i.e., the initial generation before any iterative correction); later steps indicate that the wrapper's internal loop produced at least one intervening response before the error was committed or persisted.

\begin{table*}[t]
\centering
\caption{Failure-step distribution from binary-search attribution on retained GAIA failed instances (GPT-4.1 backbone). Step index is the earliest conversation turn at which the agent output diverges from the ground truth. Jarvis (88 analyzed instances) commits errors overwhelmingly at step~1, consistent with its single-shot handoff. ChatEval (101 analyzed instances) also concentrates errors at step~1, showing that the debate loop does not reliably rescue initial-turn failures; a smaller fraction persists or is first introduced at later odd-numbered debate rounds.}
\label{tab:appendix_failure_step_distribution}
\small
\setlength{\tabcolsep}{5pt}
\renewcommand{\arraystretch}{1.12}
\resizebox{\textwidth}{!}{
\begin{tabular}{lrrrrrrr}
    \toprule
    \textbf{Wrapper} & \textbf{Step 1} & \textbf{Step 3} & \textbf{Step 5} & \textbf{Step 7} & \textbf{Step 11} & \textbf{Step 13} & \textbf{Total analyzed} \\
    \midrule
    Jarvis  & 72 (81.8\%) & --- & --- & --- & --- & --- & 88 \\
    ChatEval & 58 (57.4\%) & 6 (5.9\%) & 2 (2.0\%) & 6 (5.9\%) & 1 (1.0\%) & 1 (1.0\%) & 101 \\
    \midrule
    \multicolumn{8}{l}{\footnotesize\textit{Note: Jarvis step=0 (5 instances, 5.7\%) and ChatEval unattributed instances are omitted. Step indices correspond to BenchAgent message-log positions.}} \\
    \bottomrule
\end{tabular}
}
\end{table*}

The failure-step profile is consistent with two complementary observations. First, both wrappers commit the majority of analyzed errors at the first agent turn, suggesting that the tested wrapper structures often do not compensate for an initial-turn knowledge or retrieval gap. Second, while ChatEval's debate loop does occasionally defer or surface an error at a later round (steps~3--13 account for~16.8\% of ChatEval failures), this correction window is limited and does not translate into a net accuracy gain on GAIA relative to the simpler Jarvis handoff. This supports the bounded process reading in Table~\ref{tab:appendix_process_diagnostics}: in these retained runs, ChatEval's additional coordination expands token usage without a proportionate observed error-correction benefit.

\subsection{Runtime-Generated Workflow Evidence}
\label{subsec:appendix_runtime_evidence}

Table~\ref{tab:appendix_claude_trj_stats} reports structured-event statistics from the retained Claude Code dumps that preserve explicit team-creation and role-message structure. ``Subagents'' counts distinct non-lead spawned agents, ``Max Depth'' is the longest directed delegation chain from the team lead, ``Verifier Triggers'' counts unique task-level non-idle verifier emissions, and ``Explicit Repair Loops'' counts verifier events that route control back to the solver or request correction.

\begin{table*}[t]
    \centering
    \caption{Structured-event statistics from retained Claude Code GAIA trajectory dumps with explicit team-creation and role-message structure.}
    \label{tab:appendix_claude_trj_stats}
    \small
    \setlength{\tabcolsep}{6pt}
    \renewcommand{\arraystretch}{1.15}
    \resizebox{\textwidth}{!}{
    \begin{tabular}{lcccc}
        \toprule
        \textbf{Trace Dump} & \textbf{Subagents} & \textbf{Max Depth} & \textbf{Verifier Triggers} & \textbf{Explicit Repair Loops} \\
        \midrule
        3.json & 5 & 1 & 0 & 0 \\
        6.json & 4 & 4 & 3 & 0 \\
        claude\_code\_logs\_gaia1.json & 12 & 2 & 0 & 0 \\
        claude\_code\_logs\_gaia3.json & 4 & 3 & 2 & 0 \\
        \midrule
        Team-run average & 6.25 & 2.50 & 1.25 & 0.00 \\
        \bottomrule
    \end{tabular}
    }
\end{table*}

These statistics support the narrower process claim made in the main text: the retained runtime-generated workflow traces repeatedly instantiate explicit subagents, reach delegation depth up to 4, and expose verifier-stage events in two structured runs. They do not support an exhaustive count of all recovery loops.

Table~\ref{tab:appendix_team_init_trace} extracts the team-initialization sequence from \texttt{6.json} item~17 (the first item with all four roles instantiated). The trace shows the team lead deciding at runtime how many agents to create, what roles to assign, and what prompt each role receives---a workflow structure that does not exist as a static template but is constructed during the GAIA run itself.

\begin{table*}[t]
    \centering
    \caption{Team-initialization trace from \texttt{6.json} item~17 (\texttt{gaia-validation-l2p15}). The team lead (CC-workflow orchestrator) emits one \texttt{TeamCreate} call followed by four \texttt{Agent} spawns in a single conversation turn, each with a distinct role prompt baked in at spawn time. Step numbers correspond to assistant message indices in the raw log.}
    \label{tab:appendix_team_init_trace}
    \footnotesize
    \setlength{\tabcolsep}{3.5pt}
    \renewcommand{\arraystretch}{1.12}
    \resizebox{\textwidth}{!}{
    \begin{tabular}{cp{0.10\textwidth}p{0.14\textwidth}p{0.58\textwidth}}
        \toprule
        \textbf{Step} & \textbf{Actor} & \textbf{Tool / Event} & \textbf{Retained excerpt (parameters condensed)} \\
        \midrule
        1 & Team Lead & \texttt{TeamCreate} & \texttt{agent\_type}: planner; \texttt{team\_name}: gaia-validation-l2p15; \textit{description}: ``Coordinated agent team for GAIA validation pipeline on \texttt{level2\_part15.jsonl}. Each role operates with bash for file I/O; outputs only to \texttt{outputs/} and \texttt{logs/}.'' \\
        \midrule
        & Runtime & spawn ack & \texttt{agent\_id}: planner@gaia-validation-l2p15; team registered; agent awaiting mailbox. \\
        \midrule
        2 & Team Lead & \texttt{Agent} (planner) & \texttt{model}: sonnet; \texttt{isolation}: worktree; \textit{prompt}: ``Use bash tools ONLY (\texttt{head}, \texttt{sed}, \texttt{tail}) to read the first~3 lines from \texttt{level2\_part15.jsonl}. For each, extract \texttt{task\_id} and question, then pass them along with their line number (0-indexed) to solver, one at a time.'' \\
        \midrule
        3 & Team Lead & \texttt{Agent} (solver) & \texttt{model}: sonnet; \textit{prompt}: ``You are the Solver. For each GAIA problem sent to you (\texttt{task\_id}, question, line number), solve based strictly on the question and associated context. Prepare a concise answer in the required GAIA format. Send your answer, \texttt{task\_id}, and reasoning to the Verifier.'' \\
        \midrule
        4 & Team Lead & \texttt{Agent} (verifier) & \texttt{model}: sonnet; \textit{prompt}: ``For each problem: receive \texttt{task\_id}, answer, and reasoning from Solver. Check GAIA answer format; assess answer validity. If correct and in format, forward to Writer with \texttt{task\_id}, answer, reasoning. If not, return to Solver noting the issue.'' \\
        \midrule
        5 & Team Lead & \texttt{Agent} (writer) & \texttt{model}: sonnet; \textit{prompt}: ``Upon receiving \texttt{task\_id}, answer, and reasoning from the Verifier, check if \texttt{outputs/<task\_id>.txt} exists. If it does, SKIP. If not, create \texttt{outputs/<task\_id>.txt} and write ONLY the answer there. Always write the full reasoning log to \texttt{logs/<task\_id>.log}.'' \\
        \midrule
        & Runtime & 4$\times$ spawn ack & \texttt{agent\_id}: solver@, verifier@, writer@gaia-validation-l2p15; all four roles live in shared mailbox. \\
        \midrule
        6 & Team Lead & \texttt{Bash} & \texttt{head -n 3 split\_validation/level2\_part15.jsonl}; first returned \texttt{task\_id}: \texttt{48eb8242{\ldots}} (``How many nonindigenous crocodiles in Florida 2000--2020?''). \\
        \bottomrule
    \end{tabular}
    }
\end{table*}

The team-initialization sequence illustrates two key properties of the runtime-generated workflow. First, the workflow structure (4-role pipeline with planner $\to$ solver $\to$ verifier $\to$ writer routing) is not hard-coded in any configuration file but is emitted live by the team lead in a single conversation turn, conditioned on the GAIA task batch. Second, each role's prompt is embedded at spawn time and specializes the subagent's behavior from the first message it receives---the verifier, for instance, is instructed from the outset to return failing answers to the solver rather than simply forward everything to the writer. This is structurally distinct from fixed MAS designs in which role assignments and routing are predetermined before any task is seen.

Table~\ref{tab:appendix_task_exec_trace} continues the trace from team initialization to full execution, following one GAIA instance (\texttt{task\_id}: \texttt{48eb8242}, ``How many nonindigenous crocodiles were found in Florida from the year 2000 through 2020?'', ground truth: 11) through each pipeline stage. The trace is drawn from \texttt{6.json} items~15, 22, 24, and~30, corresponding to the Planner, Solver, Verifier, and Writer turns respectively for this instance.

\begin{table*}[t]
    \centering
    \caption{Task-execution trace for CC-workflow instance \texttt{48eb8242} (GAIA Level~2, Florida crocodile count task, ground truth: 11). Steps E1--E11 span four role agents across items~15, 22, 24, and~30 of \texttt{6.json}. Planner dispatches via a \texttt{SendMessage}; Solver resolves and forwards; Verifier checks format and routes; Writer enforces the skip-if-exists guard before committing. All role-level prompts were set at spawn time (Table~\ref{tab:appendix_team_init_trace}).}
    \label{tab:appendix_task_exec_trace}
    \footnotesize
    \setlength{\tabcolsep}{3.5pt}
    \renewcommand{\arraystretch}{1.12}
    \resizebox{\textwidth}{!}{
    \begin{tabular}{cp{0.09\textwidth}p{0.14\textwidth}p{0.60\textwidth}}
        \toprule
        \textbf{Step} & \textbf{Role} & \textbf{Tool / Event} & \textbf{Retained excerpt} \\
        \midrule
        E1 & Planner & \texttt{Bash} & \texttt{head -n 3 split\_validation/level2\_part15.jsonl}; returns \texttt{task\_id}: \texttt{48eb8242{\ldots}}, \textit{Question}: ``How many nonindigenous crocodiles were found in Florida from the year 2000 through 2020?~You can get the data from the USGS Nonindigenous Aquatic Species database.'' \\
        \midrule
        E2 & Planner & \texttt{SendMessage} (to solver) & Dispatches task to Solver's mailbox: ``Level2\_part15.jsonl line~0: task\_id=48eb8242-{\ldots}, question=How many nonindigenous crocodiles\ldots''; Solver receives as a teammate-message. \\
        \midrule
        E3 & Solver & [teammate-message from Planner] & Receives dispatch payload containing \texttt{task\_id} and full question text; triggers Solver's internal answer generation. \\
        \midrule
        E4 & Solver & [internal generation] & ``Answer: 11. Reasoning: According to the USGS Nonindigenous Aquatic Species database, there were 11 confirmed occurrences of nonindigenous crocodiles in Florida between 2000 and 2020.'' \\
        \midrule
        E5 & Solver & \texttt{SendMessage} (to Verifier) & \textit{message}: ``Answer: 11 / Task ID: 48eb8242{\ldots} / Reasoning: USGS NAS database, 11 confirmed occurrences of nonindigenous Crocodylus in Florida 2000--2020.''; routing: sender=solver (green), target=@Verifier. \\
        \midrule
        E6 & Verifier & [teammate-message from Solver] & Receives \texttt{task\_id}, answer, and reasoning; role instruction: ``Check Gaia answer format; if correct and in format, forward to Writer; if not, return to Solver noting the issue.'' \\
        \midrule
        E7 & Verifier & [format/validity check] & Answer ``11'' is a valid integer in GAIA format; no extra fields; reasoning is consistent. Proceeds to forward. \\
        \midrule
        E8 & Verifier & \texttt{SendMessage} (to writer) & \textit{summary}: ``Verified: answer and format for task 48eb8242{\ldots}''; \textit{content}: ``Task ID: 48eb8242{\ldots} / Answer: 11 / Reasoning: \ldots''; routing: sender=verifier (yellow), target=@writer (purple). \\
        \midrule
        E9 & Writer & [teammate-message from Verifier] & Receives: ``Task ID: 48eb8242{\ldots} / Answer: 11 / Reasoning: USGS NAS{\ldots}''; role instruction: ``check if \texttt{outputs/<task\_id>.txt} exists; if so, SKIP; if not, write ONLY the answer.'' \\
        \midrule
        E10 & Writer & \texttt{Read} (existence guard) & \texttt{file\_path}: \texttt{outputs/48eb8242{\ldots}.txt}; checks whether a prior write already committed this answer (skip-if-exists gate). \\
        \midrule
        E11 & Writer & \texttt{Write} (log commit) & \texttt{file\_path}: \texttt{logs/48eb8242{\ldots}.log}; writes full reasoning log; answer ``11'' written to \texttt{outputs/48eb8242{\ldots}.txt}. Task finalized. \\
        \bottomrule
    \end{tabular}
    }
\end{table*}

The execution trace exposes three runtime behaviors that are not present as logged state transitions in the fixed MAS instantiations evaluated here. First, inter-role communication is explicit and visible: each handoff from Planner $\to$ Solver $\to$ Verifier $\to$ Writer appears as a \texttt{SendMessage} event with a typed payload, making the state at each boundary auditable from the retained logs. Second, the verifier applies a format gate at step E7, distinguishing answer delivery from a simple pass-through chain; under the role prompt, a failed check would route control back to the Solver. Third, the Writer applies a skip-if-exists guard at step E10, preventing double-writes even if the same verified answer arrives via a separate message path. These state transitions are the trace-visible correlates of the candidate mechanisms in Table~\ref{tab:gaia_mechanism_evidence}; they support the existence of verifier-stage control and explicit state handoff in this retained execution, while leaving their causal contribution to accuracy for ablation.

Table~\ref{tab:appendix_cc_workflow_role_prompts} records the full spawn-time role prompts for each of the four agents. These prompts are extracted verbatim from the \texttt{prompt} parameter of each \texttt{Agent} call in \texttt{6.json} item~17 (steps~2--5 of Table~\ref{tab:appendix_team_init_trace}). Because the prompts are issued at spawn time and not updated within a run, they represent the complete behavioral contract for each role throughout the session: the only subsequent instructions each agent receives come from the shared mailbox.

\begin{table*}[t]
\centering
\caption{CC-workflow runtime spawn-time role prompts extracted verbatim from \texttt{6.json} item~17. Each prompt constitutes the full behavioral contract for the corresponding role; no further system-level instruction is issued after spawn. The prompts show role specialization (file-read restriction for Planner, format commitment for Solver, routing logic for Verifier, idempotent write guard for Writer) arising from a single live decision by the team lead, not from a pre-written template.}
\label{tab:appendix_cc_workflow_role_prompts}
\footnotesize
\setlength{\tabcolsep}{4pt}
\renewcommand{\arraystretch}{1.12}
\begin{tabular}{p{0.10\textwidth}p{0.85\textwidth}}
\toprule
\textbf{Role} & \textbf{Spawn-time prompt (verbatim)} \\
\midrule
Planner & Use bash tools ONLY (head, sed, tail, etc.) to read the first 3 lines from \texttt{split\_validation/level2\_part15.jsonl}. For each, extract the \texttt{task\_id} and question, then pass them along with their line number (0-indexed) to solver, one at a time. Output nothing else. Do not use internal file reading functions. \\
\midrule
Solver & You are the Solver. For each GAIA problem sent to you (\texttt{task\_id}, question, line number), solve based strictly on the question and associated context. Prepare a concise answer string in the required GAIA answer format. Send your answer, \texttt{task\_id}, and reasoning to the Verifier. Output nothing else. \\
\midrule
Verifier & For each problem: receive \texttt{task\_id}, answer, and reasoning from Solver. Check GAIA answer format (no extra data, correct fields, etc.); assess answer validity as best as possible. If correct and in format, forward to Writer with \texttt{task\_id}, answer, reasoning. If not, return to Solver noting the issue. Output nothing else. \\
\midrule
Writer & Upon receiving a \texttt{task\_id}, answer, and reasoning from the verifier, check if \texttt{outputs/<task\_id>.txt} exists. If it does, SKIP. If not, create \texttt{outputs/<task\_id>.txt} and write ONLY the answer there, with no extra output. Always also write the full reasoning log to \texttt{logs/<task\_id>.log}. Never overwrite existing files. Output nothing else. \\
\bottomrule
\end{tabular}
\end{table*}

Comparing Table~\ref{tab:appendix_cc_workflow_role_prompts} with the reproduced fixed MAS prompt contracts in Table~\ref{tab:appendix_reproduced_mas_contracts} reveals a structural difference in how behavioral constraints are specified. Fixed MAS prompts (Jarvis, ChatEval, LLM-Debate, CAMEL) are written as general-purpose expert personas: they describe a type of reasoning (mathematical, logical, critical) that applies to any task, and the role assignment does not change across instances. The CC-workflow prompts, by contrast, are operationally specific to the current batch: the Planner prompt names the exact JSONL file (\texttt{level2\_part15.jsonl}) and prescribes the exact tools (\texttt{head}, \texttt{sed}, \texttt{tail}); the Solver prompt specifies the exact target for routing (``send answer, \texttt{task\_id}, and reasoning to the Verifier''); the Verifier prompt specifies an explicit fallback path (``if not, return to Solver noting the issue''); and the Writer prompt encodes an idempotency contract (``if \texttt{outputs/<task\_id>.txt} exists, SKIP''). This specificity is only possible because the team lead generates these prompts in real time, after reading the task batch and before any agent begins work.

A second distinguishing feature is the scope of the behavioral contract. In fixed MAS designs, the system-level prompt defines the full behavioral contract; agents may receive task-specific content via the conversation, but their behavioral role does not change. In the CC-workflow, the spawn-time prompt is the full behavioral contract for each role: there is no separate system-level prompt, and the only post-spawn input each agent receives is the mailbox traffic from other roles. The Solver in Table~\ref{tab:appendix_cc_workflow_role_prompts}, for instance, is aware from its first token that it must route its output to the Verifier specifically---it does not receive this instruction as part of the task, and it cannot deviate from it without violating the prompt it was initialized with. This is the mechanism by which role specialization persists throughout a run without a central dispatcher repeatedly re-stating the routing rules.

\begin{figure*}[t]
    \centering
    \includegraphics[width=0.85\textwidth]{\detokenize{resources/figure/same-task workflow comparison.pdf}}
    \caption{Controlled same-task workflow comparison between a fixed MAS pipeline and a runtime-generated workflow. The fixed MAS passes forward only a partial intermediate result and loses the instruction-level constraint, producing a locally correct computation but a globally incorrect final answer. The retained Claude Code trace preserves context and reasoning in the state payload and includes verifier-stage evidence before final answer release.}
    \label{fig:appendix_same_task_workflow_comparison}
\end{figure*}

\begin{figure*}[t]
    \centering
    \includegraphics[width=0.85\textwidth]{resources/figure/evo_true_claude_true.pdf}
    \caption{Supplementary matched-success GAIA case. Both EvoAgent and Claude Code eventually solve the task, but Claude Code externalizes intermediate state and routes it through verification more explicitly.}
    \label{fig:appendix_gaia_generated_team_success}
\end{figure*}

\paragraph{Scope and limitations of the retained trace evidence.}
The trace evidence collected in Tables~\ref{tab:appendix_team_init_trace}--\ref{tab:appendix_cc_workflow_role_prompts} and in the preceding trajectory excerpts supports a narrow set of claims. First, the team-initialization trace confirms that the CC-workflow creates its agent structure dynamically and assigns role-specific prompts at the point of instantiation, not from a fixed template. Second, the task-execution trace confirms that role-to-role handoffs are mediated by explicit \texttt{SendMessage} events, that the verifier stage receives the answer before the writer stage receives it, and that the writer enforces an idempotent write guard. Third, the spawn-time prompt table confirms that behavioral specialization for this particular GAIA batch is embedded in the initial prompt for each agent, not relayed through subsequent messages.

The retained traces are intentionally used for bounded process claims. The detailed tables are drawn from one GAIA batch partition (\texttt{level2\_part15.jsonl}), so they establish that this runtime-generated structure occurred, not that the same four-role structure is universal across all CC-workflow batches. Because the team lead generates structure conditioned on the task batch, different inputs could yield different role counts or routing graphs. Trace visibility is also partial: spawn-time prompts and inter-role \texttt{SendMessage} events are visible, while provider-side bookkeeping such as context compaction, token caching, and background retry logic is outside the retained artifacts. The ``explicit repair loop'' count in Table~\ref{tab:appendix_claude_trj_stats} is therefore a lower-bound event count, not a claim that no unobserved repair behavior occurred.

The evidence tables therefore support the descriptive claim in the main text: the retained CC-workflow traces instantiate explicit subagents, implement a verifier-stage gate, and preserve routing state through visible handoffs. They do not assign the accuracy advantage to one mechanism. In particular, the logs alone cannot separate the effects of the verifier gate, the skip-if-exists idempotency guard, spawn-time prompt specificity, context packaging, and other runtime behaviors. Controlled ablations that disable individual design choices while holding all others fixed would be required for that causal attribution, as noted in Table~\ref{tab:gaia_mechanism_evidence}.

\subsection{Responsible Research and Artifact Use}
\label{subsec:appendix_responsible_artifact_use}

\paragraph{Artifact use.}
This paper uses public benchmarks and baseline descriptions for research evaluation, with dataset sampling, tool regimes, protocol alignment, and retained trace evidence documented in Appendices~\ref{subsec:appendix_protocol_fidelity}--\ref{subsec:appendix_runtime_evidence}. We follow the original license and access terms for the benchmarks, baselines, and tool artifacts used in the study, and we do not redistribute restricted source data, credentials, private paths, or provider tokens.

\paragraph{Release plan.}
During anonymous review, we do not publicly release non-anonymized code, local configuration, or raw traces that could expose private infrastructure, credentials, local paths, or non-public runtime details. Upon acceptance, we plan to release the evaluation code, analysis scripts, anonymized or non-sensitive aggregate statistics, and trajectory tooling needed to reproduce the reported tables and audit the retained process evidence where redistribution is permitted.

\paragraph{Risks and AI assistance.}
The main empirical risks are the one-run pass@1 setting, the PAE nature of the CC-workflow comparison, and the possibility that agent-count summaries obscure workflow, tool-surface, and context-management effects; these risks are discussed in the main-text limitations and mechanism-evidence notes. AI assistants were used for language editing, code assistance, and drafting support. The authors reviewed and take responsibility for the final claims, evidence, and manuscript text.

\subsection{Additional Broad-Benchmark Analyses}
\label{subsec:appendix_broad_benchmark_analyses}

Appendix Figures~\ref{fig:appendix_main_exp1_summary_bars}--\ref{fig:appendix_delta_lollipops} collect the supplementary broad-benchmark views. Figure~\ref{fig:appendix_main_exp1_summary_bars} summarizes the three aggregate rows used in the main-text cost discussion, while Figures~\ref{fig:appendix_accuracy_heatmap}--\ref{fig:appendix_delta_lollipops} show task-level accuracy, workflow lift, per-benchmark rankings, and per-workflow benchmark-level lift profiles.

\begin{figure*}[htbp]
    \centering
    \includegraphics[width=0.85\textwidth]{resources/figure/main_experiment1_overview/method_summary_bars.png}
    \caption{Main Experiment~1 method-level summaries. The panels visualize the summary rows of Table~\ref{tab:main_exp1_overall}: benchmark-balanced average accuracy, instance-level average end-to-end token usage, and instance-level average execution time.}
    \label{fig:appendix_main_exp1_summary_bars}
\end{figure*}

\begin{figure*}[htbp]
    \centering
    \includegraphics[width=0.85\textwidth]{resources/figure/main_experiment1_overview/accuracy_heatmap.png}
    \caption{Benchmark-by-method accuracy heatmap for Main Experiment~1. Each cell reports benchmark-level pass@1 accuracy for one workflow on one benchmark.}
    \label{fig:appendix_accuracy_heatmap}
\end{figure*}

\begin{figure*}[htbp]
    \centering
    \includegraphics[width=0.85\textwidth]{resources/figure/main_experiment1_overview/workflow_lift_heatmap.png}
    \caption{Workflow lift relative to the Single Agent baseline in Main Experiment~1. Green cells indicate positive accuracy delta, and orange/red cells indicate lower accuracy than the baseline on the same benchmark.}
    \label{fig:appendix_delta_heatmap}
\end{figure*}

\begin{figure*}[htbp]
    \centering
    \includegraphics[width=0.85\textwidth]{resources/figure/main_experiment1_overview/benchmark_dotpanels.png}
    \caption{Per-benchmark method comparison for Main Experiment~1. Each panel shows the pass@1 accuracy distribution across workflows for one benchmark, making task-specific outliers and near-ties easier to inspect than in the wide table.}
    \label{fig:appendix_benchmark_dotpanels}
\end{figure*}

\begin{figure*}[htbp]
    \centering
    \includegraphics[width=0.85\textwidth]{resources/figure/main_experiment1_overview/workflow_lift_lollipops.png}
    \caption{Per-workflow benchmark-level lift profiles relative to the Single Agent baseline. Points to the right of zero indicate positive benchmark-level accuracy delta, while points to the left indicate lower accuracy than the matched baseline.}
    \label{fig:appendix_delta_lollipops}
\end{figure*}
